\documentclass[aspectratio=169]{beamer}
\usepackage[T1]{fontenc}
\usepackage{tikz}
\usetikzlibrary{arrows.meta}
\setbeamertemplate{navigation symbols}{}
\definecolor{papel}{HTML}{F7F5EF}
\definecolor{tinta}{HTML}{192A34}
\definecolor{gris}{HTML}{4B5C65}
\definecolor{verde}{HTML}{00665B}
\definecolor{ocre}{HTML}{934219}
\definecolor{regla}{HTML}{CDD2CB}
\definecolor{codigo}{HTML}{EAEFEA}
\setbeamercolor{background canvas}{bg=papel}

% Coordenadas compartidas en centímetros para una página Beamer de 16 × 9.
\newenvironment{lienzo}{%
  \begin{tikzpicture}[remember picture,overlay,x=1cm,y=-1cm]
  \begin{scope}[shift={(current page.north west)}]
  \tikzset{
    txt/.style={anchor=north west,inner sep=0pt,text=tinta,align=left},
    mono/.style={txt,font=\ttfamily},
    flecha/.style={draw=verde,line width=.85pt,-{Stealth[length=2.1mm]}}
  }
}{%
  \end{scope}
  \end{tikzpicture}
}
\newcommand{\cabecera}[3]{%
  \fill[#1] (.85,.62) rectangle (.97,1);
  \node[txt,text=#1,font=\bfseries\small] at (1.12,.68) {#2};
  \node[txt,font=\bfseries\LARGE] at (.85,1.27) {#3};
  \draw[regla] (.85,2.38)--(15.15,2.38);
}
\newcommand{\pie}[1]{%
  \draw[regla] (.85,8.18)--(15.15,8.18);
  \node[txt,text=gris,font=\small] at (.85,8.45) {#1};
  \node[txt,text=gris,font=\small,anchor=north east] at (15.15,8.45) {\insertframenumber};
}
\newcommand{\bloque}[5]{%
  \node[txt,text=#1,font=\bfseries\large] at (#2,#3) {#4};
  \node[txt,text width=6.55cm,font=\normalsize] at (#2,{#3+.60}) {#5};
}
\newcommand{\etiqueta}[4]{%
  \node[txt,text=#1,font=\small\bfseries] at (#2,#3) {#4};
}
\begin{document}

\begin{frame}[plain,noframenumbering]
\begin{lienzo}
\fill[tinta] (0,0) rectangle (16,9);
\fill[verde] (.94,1.13) rectangle (2.53,1.27);
\fill[ocre] (2.63,1.13) rectangle (3.55,1.27);
\node[txt,text=white,font=\small\bfseries] at (.94,1.65) {EXPOSICIÓN DEL PRIMER CORTE · GA-SUM-02 / EXP-1};
\node[txt,text=white,font=\bfseries\fontsize{29}{34}\selectfont] at (.90,2.45) {Servicios web};
\node[txt,text=white,font=\bfseries\LARGE] at (.94,3.87) {SOAP y REST};
\node[txt,text=white,font=\large] at (.94,4.69) {Análisis de un sistema distribuido real};
\draw[white,opacity=.4] (.94,5.86)--(15.06,5.86);
\node[txt,text=white,font=\small\bfseries] at (.94,6.15) {Integrantes · Grupo CCEG};
\node[txt,text=white,font=\small] at (.94,6.65) {Calderon Saltos Joseph Alexander\\
Coronel Gomez Khriz Alexander\\
Emanuel Pino Juliana Romina\\
Garcia Bazurto Kevin Leonel};
\node[txt,text=white,font=\small\bfseries] at (9.2,6.15) {Docente};
\node[txt,text=white,font=\small] at (9.2,6.65) {Guerrero Ulloa Gleiston Ciceron};
\node[txt,text=white,font=\small\bfseries] at (9.2,7.24) {Aplicaciones Distribuidas};
\node[txt,text=white,font=\small] at (9.2,7.71) {7.º nivel · Paralelo A};
\node[txt,text=white,font=\small] at (9.2,8.31) {16 de noviembre de 2026};
\end{lienzo}
\end{frame}



\begin{frame}[plain]
\begin{lienzo}
\cabecera{verde}{TEORÍA}{Fundamentos de los servicios web y SOAP}
\node[txt,text width=14cm,font=\large] at (.85,2.75) {Un servicio web permite que aplicaciones con tecnologías distintas intercambien información a través de una red.};
\draw[flecha] (3.82,4.27)--(12.02,4.27);
\draw[flecha] (12.02,5.06)--(3.82,5.06);
\node[txt,font=\bfseries\Large] at (.85,4.41) {Cliente};
\node[txt,font=\bfseries\Large,anchor=north east] at (15.15,4.41) {Servidor};
\node[txt,anchor=south,font=\small] at (8,4.15) {Solicitud};
\node[txt,anchor=north,font=\small] at (8,5.18) {Respuesta};
\bloque{ocre}{.85}{6.05}{SOAP · Protocolo de mensajería}{Intercambia mensajes estructurados en XML dentro de un sobre (\textit{Envelope}).}
\bloque{ocre}{8.45}{6.05}{WSDL · Contrato del servicio}{Describe qué operaciones ofrece el servicio y cómo acceder a ellas.}
\pie{Fuentes · W3C, SOAP 1.2 y WSDL 2.0}
\end{lienzo}
\end{frame}

\begin{frame}[plain]
\begin{lienzo}
\cabecera{verde}{TEORÍA}{REST y sus diferencias con SOAP}
\node[txt,text width=14cm,font=\large] at (.85,2.76) {REST es un estilo arquitectónico para sistemas distribuidos.};
\draw[regla] (8,3.55)--(8,7.60);
\bloque{verde}{.85}{3.61}{Los recursos en REST}{La URI identifica el recurso. HTTP ofrece métodos como GET, POST, PUT y DELETE.}
\bloque{verde}{.85}{6.03}{Cada solicitud trae su contexto}{REST separa cliente y servidor. Cada solicitud lleva la información necesaria para procesarla.}
\bloque{ocre}{8.45}{3.61}{Mensajes y formatos}{SOAP usa mensajes XML. REST admite representaciones como JSON y XML.}
\bloque{ocre}{8.45}{6.03}{Contrato y acceso}{WSDL es habitual en SOAP. En REST, el método HTTP y la URI indican qué se pide y sobre qué recurso.}
\pie{Fuentes · Fielding (2000), REST. IETF, RFC 9110 (2022).}
\end{lienzo}
\end{frame}

\begin{frame}[plain]
\begin{lienzo}
\cabecera{verde}{EJERCICIO EN POSTMAN}{REST lee el ID de la URL}
\fill[codigo] (.85,2.88) rectangle (7.55,6.91);
\fill[codigo] (8.06,2.88) rectangle (15.15,6.91);
\etiqueta{verde}{1.13}{3.2}{SOLICITUD · SIN CUERPO}
\etiqueta{verde}{8.35}{3.2}{RESPUESTA · 200 OK · JSON}
\node[mono] at (1.13,3.92) {GET /api/practica/ordenes/\textcolor{verde}{\textbf{2}}};
\node[mono] at (8.35,3.92) {\{\\
\quad "idOrden": 2,\\
\quad "estadoOrden": "PENDIENTE",\\
\quad "total": 47.50\\
\}};
\node[txt,text width=13.9cm] at (.85,7.24) {\textbf{En Java} \enspace \texttt{@PathVariable} toma el identificador de la ruta.};
\pie{Orden 2 de papelería · JSON sencillo de consumir desde el frontend}
\end{lienzo}
\end{frame}

\begin{frame}[plain]
\begin{lienzo}
\cabecera{ocre}{EJERCICIO EN POSTMAN}{SOAP lee el ID del cuerpo XML}
\fill[codigo] (.85,2.88) rectangle (7.55,6.15);
\fill[codigo] (8.06,2.88) rectangle (15.15,6.15);
\etiqueta{ocre}{1.13}{3.17}{SOLICITUD · POST /ws}
\etiqueta{ocre}{8.35}{3.17}{RESPUESTA · 200 OK}
\node[mono,font=\ttfamily\small] at (1.13,3.78) {\textless oc:ConsultarOrdenRequest\textgreater\\
\quad\textless oc:idOrden\textgreater\textcolor{ocre}{\textbf{2}}\textless/oc:idOrden\textgreater\\
\textless/oc:ConsultarOrdenRequest\textgreater};
\node[mono,font=\ttfamily\small] at (8.35,3.78) {\textless idOrden\textgreater2\textless/idOrden\textgreater\\
\textless estadoOrden\textgreater PENDIENTE\textless/estadoOrden\textgreater\\
\textless total\textgreater47.50\textless/total\textgreater};
\node[txt,text width=14cm] at (.85,6.49) {\textbf{En Java} \enspace Ambos métodos llaman a \texttt{ordenCompraService.obtener(id)}.};
\node[txt,text width=14cm] at (.85,7.22) {\textbf{Decisión} \enspace REST para el frontend. SOAP si la integración necesita WSDL.};
\pie{Extractos del cuerpo de la solicitud y del cuerpo de la respuesta SOAP}
\end{lienzo}
\end{frame}

\begin{frame}[plain]
\begin{lienzo}
\fill[tinta] (0,0) rectangle (16,9);
\fill[verde] (.95,1.22) rectangle (2.65,1.32);
\node[txt,text=white,font=\bfseries\Huge] at (.95,2.20) {Gracias};
\node[txt,text=white,font=\large] at (.95,3.57) {SOAP y REST};
\draw[white,opacity=.35] (.95,6.15)--(15.05,6.15);
\node[txt,text=white,font=\normalsize] at (.95,6.62) {Pendientes: arquitectura y análisis del incidente};
\end{lienzo}
\end{frame}

\end{document}
